\documentclass[aps,preprint]{revtex4}%
\usepackage{amssymb}
\usepackage{amsfonts}
\usepackage{amsmath}
\usepackage{graphicx}%
\setcounter{MaxMatrixCols}{30}
%TCIDATA{OutputFilter=latex2.dll}
%TCIDATA{Version=5.50.0.2960}
%TCIDATA{CSTFile=revtex4.cst}
%TCIDATA{Created=Saturday, November 24, 2018 10:59:10}
%TCIDATA{LastRevised=Saturday, March 09, 2019 16:47:22}
%TCIDATA{<META NAME="GraphicsSave" CONTENT="32">}
%TCIDATA{<META NAME="SaveForMode" CONTENT="1">}
%TCIDATA{BibliographyScheme=Manual}
%TCIDATA{<META NAME="DocumentShell" CONTENT="Articles\SW\REVTeX 4">}
%BeginMSIPreambleData
\providecommand{\U}[1]{\protect\rule{.1in}{.1in}}
%EndMSIPreambleData
\newtheorem{theorem}{Theorem}
\newtheorem{acknowledgement}[theorem]{Acknowledgement}
\newtheorem{algorithm}[theorem]{Algorithm}
\newtheorem{axiom}[theorem]{Axiom}
\newtheorem{claim}[theorem]{Claim}
\newtheorem{conclusion}[theorem]{Conclusion}
\newtheorem{condition}[theorem]{Condition}
\newtheorem{conjecture}[theorem]{Conjecture}
\newtheorem{corollary}[theorem]{Corollary}
\newtheorem{criterion}[theorem]{Criterion}
\newtheorem{definition}[theorem]{Definition}
\newtheorem{example}[theorem]{Example}
\newtheorem{exercise}[theorem]{Exercise}
\newtheorem{lemma}[theorem]{Lemma}
\newtheorem{notation}[theorem]{Notation}
\newtheorem{problem}[theorem]{Problem}
\newtheorem{proposition}[theorem]{Proposition}
\newtheorem{remark}[theorem]{Remark}
\newtheorem{solution}[theorem]{Solution}
\newtheorem{summary}[theorem]{Summary}
\newenvironment{proof}[1][Proof]{\noindent\textbf{#1.} }{\ \rule{0.5em}{0.5em}}
\begin{document}
\preprint{HEP/123-qed}
\title[Short title for running header]{REV\TeX4}
\author{First Author}
\affiliation{My Institution}
\author{Second Author}
\affiliation{My Institution}
\author{Third Author}
\affiliation{Other Institution}
\keywords{one two three}
\pacs{PACS number}

\begin{abstract}
Shell document for REV\TeX{} 4.

\end{abstract}
\volumeyear{year}
\volumenumber{number}
\issuenumber{number}
\eid{identifier}
\date[Date text]{date}
\received[Received text]{date}

\revised[Revised text]{date}

\accepted[Accepted text]{date}

\published[Published text]{date}

\startpage{101}
\endpage{102}
\maketitle
\tableofcontents


\section{Representations of $U(2)$}%

\begin{equation}
u\left(  1,\mathbb{Q}\right)  \cong su\left(  2,\mathbb{C}\right)  \cong
so\left(  3,\mathbb{R}\right)
\end{equation}
The corresponding Lie groups are locally isomorphic and either globally
isomorphic or homomorphic images of their universal covering group. Since
$SU\left(  2,\mathbb{C}\right)  \cong U\left(  1,\mathbb{Q}\right)  $ are the
covering groups of $SO\left(  3,\mathbb{R}\right)  .$ 

The group $U\left(  2,\mathbb{C}\right)  $ is not simply connected thus has
non trivial universal covering group.

$\lceil$The Pauli matrices, $\left\{  \left.  \mathbb{\sigma}_{j}\right\vert
1\leq j\leq3\right\}  ,$ represent the cartesian operator components of a
general 3D spin operator in the spinor formulation i.e.
\begin{equation}
\Omega=\Omega_{x}\mathbb{\sigma}_{1}+\Omega_{y}\mathbb{\sigma}_{2}+\Omega
_{z}\mathbb{\sigma}_{3};\left(  \Omega_{1},\Omega_{2},\Omega_{3}\right)
\in\mathbb{R}^{3}.
\end{equation}
They also generate representations of the Lie algebras $su(2)$ and $sl\left(
2,\mathbb{C}\right)  $ i.e.
\begin{align}
su(2)  &  =\mathbb{R}\text{-}\operatorname*{linspan}\left\{  \left.
i\sigma_{j}\right\vert 1\leq j\leq3\right\} \\
sl\left(  2,\mathbb{C}\right)   &  =\mathbb{C}\text{-}\operatorname*{linspan}%
\left\{  \left.  i\sigma_{j}\right\vert 1\leq j\leq3\right\}  ,
\end{align}
which are closed under the commutator product (thus sub Lie algebras of
$M_{2}\left(  \mathbb{C}\right)  $) i.e.
\begin{align}
\left[  i\sigma_{j},i\sigma_{k}\right]   &  =-\epsilon_{jkl}i\sigma_{l},\\
\left[  i\sigma_{0},i\sigma_{j}\right]   &  =0;1\leq j\leq3,
\end{align}
where $\epsilon_{jkl}$ is the Levi-Civita symbol
\begin{equation}
\epsilon_{\sigma\left(  1,2,3\right)  }=\left(  -1\right)  ^{\pi\left(
\sigma\right)  }.
\end{equation}
They generate a representations of $SU(2)$
\begin{equation}
u\left(  \widehat{\mathbf{n}}\cdot\mathbf{\sigma}\right)  =\exp\left\{
\frac{i}{2}\left(  \widehat{\mathbf{n}}\cdot\mathbf{\sigma}\right)  \right\}
,
\end{equation}
and $SL\left(  2,\mathbb{C}\right)  $ i.e.
\begin{equation}
g\left(  \widehat{\mathbf{z}}\cdot\mathbf{\sigma}\right)  .
\end{equation}
The corresponding representation of $u(2)$ is produced by
\begin{equation}
u(2)=\operatorname*{linspan}\left\{  \left.  i\sigma_{j}\right\vert 0\leq
j\leq3\right\}  .
\end{equation}


Eulers formula can expressed in terms of the Pauli matrices as
\begin{equation}
\exp\left\{  \frac{i}{2}\left(  \widehat{\mathbf{n}}\cdot\mathbf{\sigma
}\right)  \right\}  =I_{2}\cos\left(  \theta\right)  +i\left(
\widehat{\mathbf{n}}\cdot\mathbf{\sigma}\right)  \sin\left(  \theta\right)  ,
\end{equation}
where
\begin{equation}
\exp\left\{  \frac{i}{2}\left(  \widehat{\mathbf{n}}\cdot\mathbf{\sigma
}\right)  \right\}
\end{equation}
is a counter clockwise rotation about the $\widehat{\mathbf{n}}$ axis and the
spin vector is given by
\begin{align}
\mathbf{S}  &  \equiv\mathbf{\sigma}=\left(  \left(
\begin{array}
[c]{cc}%
0 & 1\\
1 & 0
\end{array}
\right)  ,\left(
\begin{array}
[c]{cc}%
0 & -i\\
i & 0
\end{array}
\right)  ,\left(
\begin{array}
[c]{cc}%
1 & 0\\
0 & -1
\end{array}
\right)  \right) \\
\widehat{\mathbf{n}}  &  \in\mathbb{R}^{3};\,\widehat{\mathbf{n}}%
\cdot\widehat{\mathbf{n}}=1.
\end{align}
$\rfloor$%

\begin{equation}
\overset{\text{algebra}}{\sum_{1\leq j\leq3}\alpha_{j}X_{j}}\quad
\overset{\exp}{\longrightarrow}\overset{\text{group}}{\quad\exp\left\{
\sum_{1\leq j\leq3}\alpha_{j}X_{j}\right\}  \,}\equiv\left(
\widehat{\mathbf{\alpha}},\left\Vert \mathbf{\alpha}\right\Vert \right)
\end{equation}


$\lceil$%
\begin{equation}
\mathbf{u}=\left(  u_{1},u_{2}.u_{3}\right)  =i\mathbf{\sigma}=\left(  \left(
%
\begin{array}
[c]{cc}%
0 & i\\
i & 0
\end{array}
\right)  ,\left(
\begin{array}
[c]{cc}%
0 & 1\\
-1 & 0
\end{array}
\right)  ,\left(
\begin{array}
[c]{cc}%
i & 0\\
0 & -i
\end{array}
\right)  \right)
\end{equation}
$\rfloor$%

\begin{align}
\left[  u_{1},u_{2}\right]   &  =2u_{3}\\
\left[  u_{2},u_{3}\right]   &  =2u_{1}\\
\left[  u_{3},u_{1}\right]   &  =2u_{2}%
\end{align}
we can form a basis for $sl\left(  2,\mathbb{C}\right)  $ by
\begin{align}
H  &  =\frac{1}{i}u_{3}=\left(
\begin{array}
[c]{cc}%
1 & 0\\
0 & -1
\end{array}
\right) \\
X  &  =\frac{1}{2i}\left(  u_{1}-iu_{2}\right)  =\left(
\begin{array}
[c]{cc}%
0 & 1\\
0 & 0
\end{array}
\right) \\
Y  &  =\frac{1}{2i}\left(  u_{1}+iu_{2}\right)  =\left(
\begin{array}
[c]{cc}%
0 & 0\\
1 & 0
\end{array}
\right)
\end{align}%
\begin{align}
\left[  H,X\right]   &  =2X\\
\left[  H,Y\right]   &  =-2Y\\
\left[  X,Y\right]   &  =H
\end{align}%
\begin{equation}
Hv=\alpha v
\end{equation}
The eigenvalues of, $H,$ $\alpha$ are known as weights, $X$ and $Y$ are known
as step and step down operators respectively.%
\begin{align}
HXv  &  =\left(  \alpha+2\right)  Xv\\
HYv  &  =\left(  \alpha-2\right)  Yv
\end{align}%
\begin{align}
Xv  &  =0\text{ or eigenvector with eigenvalue }\left(  \alpha+2\right) \\
Yv  &  =0\text{ or eigenvector with eigenvalue }\left(  \alpha-2\right)
\end{align}
Let \ $V$ be an irreducible subspace for the Lie algebra $sl\left(
2,\mathbb{C}\right)  $ then $H$ can only have finitely many eigenvalues thus
$H$ must have a zeo eigenvlaue corresponding to the eigenvector $v_{0}$ and
one can obtain a chain of $m+1$ vectors
\begin{equation}
v_{k}=Y^{k}v_{0}%
\end{equation}
such that
\begin{align}
Yv_{m}  &  =0\\
Yv_{k}  &  =v_{k+1};0<k<m\\
Xv_{0}  &  =0\\
Xv_{k}  &  =k\left(  m-\left(  k-1\right)  \right)  v_{k-1};m>k>0\\
Hv_{k}  &  =\left(  m-2k\right)  v_{k}\\
V  &  =\mathbb{C}\text{-}\operatorname*{linspan}\left\{  v_{0},\cdots
,v_{m}\right\}
\end{align}


\begin{conclusion}
For each non-negative integer $m$, there is a unique irreducible
representation with highest weight $m$. Each irreducible representation is
equivalent to one of these. The representation with highest weight $m$ has
dimension $m+1$ with weights $m,m-2,\ldots,-(m-2),-m$ each having multiplicity one.
\end{conclusion}

\paragraph{Casimer Operator}%

\begin{align}
C  &  =-\left(  u_{1}^{2}+u_{2}^{2}+u_{3}^{2}\right) \\
&  =\left(  X+Y\right)  ^{2}-\left(  -X+Y\right)  ^{2}+H^{2}=4XY+H^{2}+2H
\end{align}%
\begin{align}
CP_{V_{m}}  &  =m\left(  m+2\right)  P_{V_{m}}\\
\left[  C,P_{V_{m}}\right]   &  =0.
\end{align}


\paragraph{Homotopy}

Let $G$ be a group with subgroup $H$ and Lie algebra $\mathfrak{g}$. Then
their is exactly one simply connected subgroup, $SG,$ with Lie algebra
$\mathfrak{g}$ and $G$ is a homomorpic image of $SG.$In fact
\begin{equation}
G\cong\left.  SG\right/  D,
\end{equation}
where $D$ is one of the discrete invariant subgroups of $SG.$
\begin{equation}
H\cong D
\end{equation}


\subsubsection{Homotopy Group}

Let $P_{1}$ and $P_{2}$ be paths in $G$ i.e paths in the parameter spaces of
the group. The the path $P_{1}P_{2}$ is the path $P_{1}$ starting and
returning to the identity followed by the path $P_{2}$ starting and returning
to the identity. Ther are as many distinct homotopy group elements as there
are classes of paths that cannot be deformed into each other.
\end{document}